% problem_id: S001-lemma-characterize-normal-crossings-normalization
% source_id: S001 (Stacks Project)
% source_file: etale.tex
% statement_locator: etale.tex lines 2566--2581
% proof_locator: etale.tex lines 2583--2676
% env: lemma
% pairing: tight (verified)
% status: complete (statement + author proof verbatim + reason + locators)
%
\begin{lemma}
\label{lemma-characterize-normal-crossings-normalization}
Let $X$ be a locally Noetherian scheme. Let $D \subset X$ be a closed
subscheme. The following are equivalent
\begin{enumerate}
\item $D$ is a normal crossings divisor in $X$,
\item $D$ is reduced, the normalization $\nu : D^\nu \to D$ is unramified,
and for any $n \geq 1$ the scheme
$$
Z_n = D^\nu \times_D \ldots \times_D D^\nu
\setminus \{(p_1, \ldots, p_n) \mid p_i = p_j\text{ for some }i\not = j\}
$$
is regular, the morphism $Z_n \to X$ is a local complete intersection
morphism whose conormal sheaf is locally free of rank $n$.
\end{enumerate}
\end{lemma}

\begin{proof}
First we explain how to think about condition (2).
The diagonal of an unramified morphism is open
(Morphisms, Lemma \ref{morphisms-lemma-diagonal-unramified-morphism}).
On the other hand $D^\nu \to D$ is separated, hence the
diagonal $D^\nu \to D^\nu \times_D D^\nu$ is closed.
Thus $Z_n$ is an open and closed subscheme of
$D^\nu \times_D \ldots \times_D D^\nu$. On the other hand,
$Z_n \to X$ is unramified as it is the composition
$$
Z_n \to D^\nu \times_D \ldots \times_D D^\nu \to \ldots \to
D^\nu \times_D D^\nu \to D^\nu \to D \to X
$$
and each of the arrows is unramified.
Since an unramified morphism is formally unramified
(More on Morphisms, Lemma
\ref{more-morphisms-lemma-unramified-formally-unramified})
we have a conormal sheaf
$\mathcal{C}_n = \mathcal{C}_{Z_n/X}$ of $Z_n \to X$, see
More on Morphisms, Definition
\ref{more-morphisms-definition-universal-thickening}.

\medskip\noindent
Formation of normalization commutes with \'etale localization by
More on Morphisms, Lemma \ref{more-morphisms-lemma-normalization-and-smooth}.
Checking that local rings are regular, or that
a morphism is unramified, or that a morphism is a
local complete intersection or that a morphism is
unramified and has a conormal sheaf which is
locally free of a given rank, may be done \'etale locally (see
More on Algebra, Lemma \ref{more-algebra-lemma-regular-etale-extension},
Descent, Lemma \ref{descent-lemma-descending-property-unramified},
More on Morphisms, Lemma \ref{more-morphisms-lemma-descending-property-lci}
and
Descent, Lemma \ref{descent-lemma-finite-locally-free-descends}).

\medskip\noindent
By the remark of the preceding paragraph and the definition
of normal crossings divisor it suffices to prove that a
strict normal crossings divisor $D = \bigcup_{i \in I} D_i$
satisfies (2). In this case $D^\nu = \coprod D_i$
and $D^\nu \to D$ is unramified (being unramified
is local on the source and $D_i \to D$ is a closed
immersion which is unramified). Similarly, $Z_1 = D^\nu \to X$
is a local complete intersection morphism because we may
check this locally on the source and each morphism $D_i \to X$
is a regular immersion as it is the inclusion of a Cartier divisor
(see Lemma \ref{lemma-strict-normal-crossings} and
More on Morphisms, Lemma \ref{more-morphisms-lemma-regular-immersion-lci}).
Since an effective Cartier divisor has an invertible
conormal sheaf, we conclude that the requirement on the
conormal sheaf is satisfied.
Similarly, the scheme $Z_n$ for $n \geq 2$ is the disjoint union
of the schemes $D_J = \bigcap_{j \in J} D_j$ where $J \subset I$
runs over the subsets of order $n$. Since $D_J \to X$ is
a regular immersion of codimension $n$
(by the definition of strict normal crossings and the
fact that we may check this on stalks by
Divisors, Lemma \ref{divisors-lemma-Noetherian-scheme-regular-ideal})
it follows in the same manner that $Z_n \to X$ has the required
properties. Some details omitted.

\medskip\noindent
Assume (2). Let $p \in D$. Since $D^\nu \to D$ is unramified, it is
finite (by Morphisms, Lemma \ref{morphisms-lemma-finite-integral}).
Hence $D^\nu \to X$ is finite unramified.
By Lemma \ref{lemma-finite-unramified-etale-local}
and \'etale localization (permissible by the discussion
in the second paragraph and the definition of normal
crossings divisors) we reduce to the case where
$D^\nu = \coprod_{i \in I} D_i$
with $I$ finite and $D_i \to U$ a closed immersion.
After shrinking $X$ if necessary, we may assume
$p \in D_i$ for all $i \in I$. The condition that $Z_1 =  D^\nu \to X$ is an
unramified local complete intersection morphism
with conormal sheaf locally free of rank $1$
implies that $D_i \subset X$ is an effective Cartier divisor, see
More on Morphisms, Lemma \ref{more-morphisms-lemma-lci} and
Divisors, Lemma \ref{divisors-lemma-regular-immersion-noetherian}.
To finish the proof we may assume $X = \Spec(A)$ is affine
and $D_i = V(f_i)$ with $f_i \in A$ a nonzerodivisor.
If $I = \{1, \ldots, r\}$, then $p \in Z_r = V(f_1, \ldots, f_r)$.
The same reference as above implies that
$(f_1, \ldots, f_r)$ is a Koszul regular ideal in $A$.
Since the conormal sheaf has rank $r$, we see that
$f_1, \ldots, f_r$ is a minimal set of generators of
the ideal defining $Z_r$ in $\mathcal{O}_{X, p}$.
This implies that $f_1, \ldots, f_r$ is a regular sequence
in $\mathcal{O}_{X, p}$ such that $\mathcal{O}_{X, p}/(f_1, \ldots, f_r)$
is regular. Thus we conclude by
Algebra, Lemma \ref{algebra-lemma-regular-mod-x}
that $f_1, \ldots, f_r$ can be extended to a regular system of parameters
in $\mathcal{O}_{X, p}$ and this finishes the proof.
\end{proof}
